\section{Discussion}
\label{sec:discussion}

% Target: 0.5 page (~400 words).
% Sources: GENERATION_FAILURE_ANALYSIS.md §Claims that need more evidence
% before publication; §Alternative hypotheses.

\subsection{What we do not claim}
\label{sec:discussion:notclaim}

% Be explicit:
% - We do not claim FTS-Diffusion as a method is fundamentally broken. The
%   stronger claim requires (i) confirming the degeneracy survives
%   sae_hidden_dim ablation (H3), (ii) seed sensitivity (H2), and (iii)
%   author confirmation that the released config is what produced the
%   paper's results.
% - We report H1 (full-budget retraining) outcome. If H1 indicates that
%   degeneracy survives the full budget, the claim is that "the released
%   configuration is structurally degenerate", not that "no configuration of
%   the architecture can work."

\subsection{Author contact}
\label{sec:discussion:contact}

% \todo{Update once author reply lands. If they confirm sae_hidden_dim=1 was
% the configuration used for the paper's headline results, that strengthens
% the reproducibility claim. If they indicate a different configuration was
% used, our finding becomes "the released artifacts do not match what was
% published — please update the release."}

\subsection{Limitations}
\label{sec:discussion:limits}

% - Single-asset focus for the most expensive experiments (H1 on S&P 500
%   only); GOOG/ZC=F results use 30-epoch checkpoints.
% - We did not run multi-seed sensitivity (H2) in time; appendix reports a
%   smaller version if available.
% - Our stress-test framework (§6) is validated on two generators; broader
%   validation across the financial-TS generator landscape is future work.
% - We focus on price-level synthetic series, mirroring the paper's setup;
%   conclusions may differ for returns-level or log-returns-level setups.
